\begin{trapbox}
\textbf{误区一：平方后的横坐标方程有实根，就一定有实交点。}
错在忽略 $B=0$ 时丢失了实纵坐标约束。
前面的竖直线三位置图已经给出反例：$x=-1$ 的平方关系有重根，但 $y^2=-4$ 无实数解。
正确做法是左右开口优先使用纵坐标二次式，或在 $B=0$ 时回原方程；此时不能用 $\Delta_x=0$ 断言相切。

\textbf{误区二：原式平方以后，随意给纵坐标选正负号。}
同一横坐标的正负纵坐标可能属于不同直线。平方消元对照图中的两条直线，横坐标方程相同，但交点纵坐标相反。
若 $B\ne0$，必须用 $y=-(Ax+C)/B$ 恢复；若 $B=0$，分别检查两种纵坐标是否存在。

\textbf{误区三：有一个交点就是切线。}
平行于轴的直线代入后降为一次关系，恰有一个交点且非相切。
只有在合法二次式存在时，重根判别才对应相切。不要跳过二次项检查。

\textbf{误区四：四种开口只需随意改一个符号。}
开口反射时，曲线和直线系数必须同步变换；交换坐标时，还要交换 $A,B$。
用前面的反射图和坐标交换图核对，不能只改公式中的 $p$。
本页约定 $p>0$，方向单独用 $s=\pm1$ 表示。例如 $y^2=4x$ 中 $p=2$，不是 $4$。

\textbf{误区五：把根差当弦长，把中点放回曲线。}
根差只是坐标投影，应由勾股关系补足距离；竖直线不能用不存在的斜率。
弦中点只须在弦所在直线上，不必在曲线上；例 2 的 $(8,17)$ 正是反例。
此外，参数候选值还须满足\textbf{两个不同实点}，例 3 的切线候选 $k=1$ 必须舍去。
\end{trapbox}
\clearpage
